\documentclass[14pt]{extarticle}

\usepackage[a4paper,margin=1in]{geometry}
\usepackage{amsmath,amssymb}
\usepackage{booktabs}
\usepackage{hyperref}

\def\mathbi#1{\textbf{\em #1}}

\title{Linear Regression}
\author{Nguyen Phu Nhan}
\date{}

\begin{document}

\maketitle

Linear regression is one of the most fundamental supervised learning algorithms. Its objective is to model the relationship between a set of input variables and a continuous target variable by fitting a linear function to observed data.

Suppose we are given a dataset consisting of $n$ observations and $p$ features. Let

\[
\mathbf{X} =
\begin{bmatrix}
x_{11} & x_{12} & \cdots & x_{1p} \\
x_{21} & x_{22} & \cdots & x_{2p} \\
\vdots & \vdots & \ddots & \vdots \\
x_{n1} & x_{n2} & \cdots & x_{np}
\end{bmatrix}
\]

be the feature matrix and

\[
\mathbi{y} =
\begin{bmatrix}
y_1 \\
y_2 \\
\vdots \\
y_n
\end{bmatrix}
\]

be the vector of target values.

The linear regression model assumes that the predicted target vector can be expressed as
\[
\hat{\mathbi{y}} = \mathbf{X}\mathbi{w} + b,
\]

where $\mathbf{X} \in \mathbb{R}^{n \times p}$ is the feature matrix containing $n$ samples and $p$ features, $\mathbi{w} \in \mathbb{R}^{p}$ is the vector of model coefficients, $b \in \mathbb{R}$ is the intercept term, and $\hat{\mathbi{y}} \in \mathbb{R}^{n}$ is the vector of predicted target values.

Our goal is to minimize the difference between the predicted target vector and the true target vector (ground truth). To quantify this difference, we introduce a loss function:
\[
L(\mathbi{w}, b)
= \frac{1}{n}
\sum_{i=1}^{n}
\frac{1}{2}
\left(\hat{y}_i - y_i\right)^2
= \frac{1}{n}
\sum_{i = 1}^n
\frac{1}{2}
\left(\mathbi{w}^\top\mathbi{x}_i + b - y_i\right)^2.
\]

When training the model, we seek the values of $\mathbi{w}$ and $b$ that minimize the loss function. This optimization problem can be written as
\[
\left(\mathbi{w}^\ast, b^\ast\right) = \operatorname*{argmin}_{\mathbi{w},\,b} L(\mathbi{w}, b).
\]

Instead of optimizing $\mathbi{w}$ and $b$ separately, we can absorb the bias
term into the feature matrix by appending a column of $1$'s to $\mathbf{X}$.
The corresponding coefficient is incorporated into $\mathbi{w}$.

The prediction function can then be written as
\[
\hat{\mathbi{y}} = \mathbf{X}\mathbi{w}.
\]

The loss function becomes
\[
L(\mathbi{w})
=
\frac{1}{2n}
\left\|
\mathbf{X}\mathbi{w}
-
\mathbi{y}
\right\|^2
=
\frac{1}{2n}
(\mathbf{X}\mathbi{w}-\mathbi{y})^\top
(\mathbf{X}\mathbi{w}-\mathbi{y}).
\]

Taking the derivative with respect to $\mathbi{w}$, we obtain
\begin{align*}
\frac{\partial L(\mathbi{w})}{\partial \mathbi{w}}
&=
\frac{\partial}{\partial \mathbi{w}}
\left[
\frac{1}{2n}
(\mathbf{X}\mathbi{w}-\mathbi{y})^\top
(\mathbf{X}\mathbi{w}-\mathbi{y})
\right] \\
&=
\frac{1}{n}
\mathbf{X}^\top
(\mathbf{X}\mathbi{w}-\mathbi{y}).
\end{align*}

Setting the derivative equal to $\mathbf{0}$ gives
\begin{align*}
\mathbf{X}^\top
(\mathbf{X}\mathbi{w}-\mathbi{y})
&=
\mathbf{0} \\
\Leftrightarrow
\mathbf{X}^\top\mathbf{X}\mathbi{w}
&=
\mathbf{X}^\top\mathbi{y}.
\end{align*}

Assuming $\mathbf{X}^\top\mathbf{X}$ is invertible, we obtain
\[
\mathbi{w} = (\mathbf{X}^\top\mathbf{X})^{-1}
\mathbf{X}^\top\mathbi{y}.
\]

\end{document}